\documentclass[a4paper]{article}
\usepackage[pdftex]{graphicx}
\usepackage[utf8]{inputenc}
\usepackage{enumerate}
\usepackage{amssymb}
\usepackage{icomma}
\usepackage{href-ul}
\hypersetup{
	colorlinks=true,
	linkcolor=blue,
	urlcolor=blue}
\usepackage{siunitx}
\sisetup{locale=DE} 
\usepackage{geometry}
\geometry{a4paper, top=15mm, left=15mm, right=15mm, bottom=15mm,
	headsep=10mm, footskip=12mm}

\begin{document}
{\bf Archimedes- Prinzip, Auftrieb in Flüssigkeiten und Gasen, Schweredruck}

\begin{enumerate}[1.]
\begin{minipage}{0.6\textwidth}
	{\bf \item Ein Kupferwürfel mit der Kantenlänge} $a=\SI{2,0}{\centi\meter}$ und der Dichte $\varrho_{Cu}=\SI{8,9}{\gram\over\centi{\meter^3}}$ wird 
	an einen Kraftmesser gehängt.  Nun wird der Würfel in einen Messkolben mit Wasser getaucht. 
	\begin{enumerate}[a)]
		\item Um wie viel ml steigt das Wasser im Kolben an ?
		\item Welche Kraft zeigt der Kraftmesser vor dem Eintauchen an? 
		\item Welche Kraft zeigt der Kraftmesser nach dem Eintauchen an?
		\item Der Versuch wird wiederholt, diesmal befindet sich der Messzylinder auf einer Waage. Sie zeigt vor dem \\ 
		Eintauchen $\SI{500}{\gram}$ an. Was zeigt sie nach dem Eintauchen an?
	\end{enumerate}
\end{minipage}
\hspace{1 cm}
\begin{minipage}{0.4\textwidth}
	\includegraphics[width=6 cm]{ach078}
\end{minipage}

{\bf \item Du möchtest zerstörungsfrei überprüfen,} ob es sich bei Deinem neuen Pokal um echtes Silber handelt.\\
Dir steht hierzu zur Verfügung:

\begin{itemize}
	\item Ein großer Meßzylinder
	\item Eine Waage
	\item Ein Brocken echtes Silber
\end{itemize}

\begin{enumerate}[a)]
	\item Welche Eigenschaften müssen der Pokal, der Meßzylinder, die Waage und der Silberbrocken haben (Größe, Beschaffenheit, ...), damit eine Überprüfung möglich ist?
	\item Erstelle eine Anleitung, wie Du bei Deiner Messung vorgehst. 
\end{enumerate}

{\bf \item In einem Heißluftballon befinden sich $\SI{5000}{\meter^3}$ Heißluft.} Sie ist um $1\over 4$ leichter als die Außenluft, welche eine Dichte von $\SI{1,3}{\kilo\gram \over {\meter^3}}$ hat.
	\begin{enumerate}[a)]
		\item Wie schwer darf die Hülle mit Korb und Besatzung höchstens sein, damit der Ballon fliegt?
		\item Welches Volumen müsste ein Heliumballon (Dichte Helium: $\varrho=\SI{178}{\gram\over{\meter^3}}$) mit gleicher Nutzlast haben? 
	\end{enumerate}


{\bf \item Natriumpolywolframat ist eine Schwerflüssigkeit } mit einer Dichte von $\varrho=\SI{3,1}{\gram\over{\centi\meter^3}}$. Auf ihr schwimmt ein Stein mit einer Masse von $\SI{80}{\gram}$. $1\over 10$ des Steins ist über der Flüssigkeitsoberfläche, $9\over 10$ darunter. Welches Volumen und welche Dichte hat der Stein? 


{\bf \item Lunerald taucht im Springerbecken (Tiefe $h=\SI{4,8}{\meter}$)bis auf den Grund}. Welcher Überdruck herrscht dort?

{\bf \item Fülle die Tabelle mit Näherungswerten aus !} Benutze hierfür die Faustregel für den Schweredruck

\renewcommand{\arraystretch}{2}
\begin{tabular}{|p{1,2 cm}|p{1,7 cm}|p{1,7 cm}|p{1,7 cm}|p{1,7 cm}|p{1,7 cm}|p{1,7 cm}|p{1,7 cm}|}
	\hline
	             & Nicht\-schwimmer\-becken & Sprung\-becken & Nordsee & Bodensee & Baikalsee & Atlantik & Mari\-anen\-graben \\
	\hline
	Tiefe h  &$\SI{1,20}{\meter}$ &$\SI{4,8}{\meter}$ &$\SI{50}{\meter}$ & &$\SI{1642}{\meter}$ & &$\SI{11}{\kilo\meter}$ \\
	\hline
	Druck & & & &$\SI{2,5}{\mega\pascal}$ & &$\SI{35}{\mega\pascal}$ & \\
	\hline
\end{tabular}

\end{enumerate} 

\fbox{
	\begin{minipage}{0.5\textwidth}
		Zur Lösung bitte \href{https://www.okuyakl.de/phys/p9achL078/ll078.pdf}{hier klicken} oder den QR-Code scannen.\\
	Weitere Arbeitsblätter gibt es unter 
	
	\href{https://www.okuyakl.de}{www.okuyakl.de}
	\end{minipage}
	\hfill
	\begin{minipage}{0.4\textwidth}
		\includegraphics[width=1.5 cm]{../../viecher/zwe03}
		\includegraphics[width=3 cm]{qr078}
		\includegraphics[width=2 cm]{../../viecher/afanticon1}
		
	\end{minipage}}
\end{document}

